\documentclass[11pt,reqno]{amsart}
\usepackage[margin=1.15in]{geometry}
\usepackage{amsmath,amssymb,amsthm,mathtools}
\usepackage{booktabs}
\usepackage[colorlinks=true,linkcolor=blue,citecolor=blue,urlcolor=blue]{hyperref}

\theoremstyle{plain}
\newtheorem{theorem}{Theorem}[section]
\newtheorem{proposition}[theorem]{Proposition}
\newtheorem{corollary}[theorem]{Corollary}
\newtheorem{conjecture}[theorem]{Conjecture}
\theoremstyle{definition}
\newtheorem{computation}[theorem]{Computation}
\newtheorem{observation}[theorem]{Observation}
\newtheorem{remark}[theorem]{Remark}

\DeclareMathOperator{\Diff}{Diff}
\DeclareMathOperator{\ord}{ord}
\DeclareMathOperator{\disc}{disc}
\DeclareMathOperator{\tr}{tr}
\DeclareMathOperator{\nrd}{nrd}
\newcommand{\Q}{\mathbb{Q}}
\newcommand{\Z}{\mathbb{Z}}
\newcommand{\C}{\mathbb{C}}
\newcommand{\HH}{\mathbb{H}}
\newcommand{\Ldual}{L^{\vee}}
\newcommand{\Lm}{L_{-}}
\newcommand{\Lp}{L_{+}}

\begin{document}

\title[The local Whittaker factor at a split level prime]
{The local Whittaker factor at a split level prime, and an incompleteness
in the standard recipe for Schofer's $\kappa$}

\author{}
\date{\today}

\begin{abstract}
Schofer's formula expresses averages of $\log\lVert\Psi_F\rVert$ over CM points on a Shimura
curve in terms of quantities $\kappa_\eta(m)$ built from derivatives of Fourier coefficients of an
incoherent Eisenstein series. The recipe used in practice to evaluate $\kappa_\eta(m)$ descends
from Kudla--Rapoport--Yang, through Schofer and Errthum, to Yang; its central step enumerates the
configurations that contribute to the derivative at the centre. We observe that this enumeration is
derived under a restriction which Kudla--Rapoport--Yang state explicitly --- that the Eisenstein
series be built from sections given by the characteristic functions of the completions of
$\mathcal{O}_k$ --- and which is not satisfied by the lattices arising for an Eichler order of
squarefree level $N$ at a \emph{split} level prime.

We compute the local Whittaker factor of the relevant $N$-scaled binary lattice at the level prime
in closed form for every valuation, and show that its value at the critical point is
$(N-1)\ord_N(m)$. We then re-run the case analysis with this factor in place and find that it
produces no additional contributing configuration, so that no $\log N$ can arise from
$\kappa_\eta(m)$ for $m>0$. Independently, a numerical evaluation of the vector-valued input form
determines, exactly, a correction of the form $A_m\log N$ that the CM values demonstrably require.
The rule producing it --- half the constant term of the input form at a nonzero isotropic coset ---
is verified against ground truth on four bases with $N\in\{2,3,5\}$, including a base with $N=5$
where the ground truth is forced by an over-determined consistency relation measured independently
of any Eisenstein computation. We identify the correction as the constant-term pairing of the input
form with the Eisenstein series attached to a nonzero isotropic coset of the discriminant form ---
an object the maximal-order literature does not compute --- and reduce it, by an oldform relation
of Schwagenscheidt, to ordinary seed-$0$ data of an enlarged lattice whose level-prime part is
unimodular; this explains structurally both the support rule and the observed normalisation. We
record the exact data any corrected formula must reproduce, and the routes we have eliminated.
\end{abstract}

\maketitle

\section{Introduction}

Let $B$ be an indefinite quaternion algebra over $\Q$ of reduced discriminant $D$, let
$\mathcal{O}\subset B$ be an Eichler order of squarefree level $N$ with $(N,D)=1$, and let
$L=\{\alpha\in\mathcal{O}:\tr\alpha=0\}$ be the associated lattice of signature $(1,2)$, so that
$\Ldual/L\cong \Z/2\times\Z/DN\times\Z/DN$ has order $2(DN)^2$ \cite[Lemma 8]{GY}. For a weakly
holomorphic vector-valued modular form $F$ of weight $1/2$ and type $\rho_L$, Borcherds' lift
$\Psi_F$ is a meromorphic modular form on $X_0^D(N)/W_{D,N}$, and Schofer's formula
\cite{Schofer,GY} evaluates
\[
  \sum_{\tau\in \mathrm{CM}(d)}\log\lvert\psi_F(\tau)\rvert
  \;=\;-\frac{\lvert \mathrm{CM}(d)\rvert}{4}\sum_{\eta\in\Ldual/L}\ \sum_{m\ge 0}
     c_\eta(-m)\,\kappa_\eta(m),
\]
valid under the hypotheses $O^+_{L,F}=O^+_L$, $c_0(0)=0$ and $c_\eta(m)\in\Z$ for $m\le 0$.

The quantities $\kappa_\eta(m)$ are not given in \cite{GY}; the reader is referred to
\cite{Errthum} and \cite{Yang} for their evaluation. Both go back to the computation of
Kudla--Rapoport--Yang \cite{KRY}, whose Proposition 2.8 identifies the configurations contributing
to the derivative of an incoherent Eisenstein series at the centre and asserts that
\emph{these are the only cases which contribute}.

\subsection*{What this note contains}
Section~\ref{sec:chain} traces that assertion to its stated scope. Section~\ref{sec:closed}
computes the missing local ingredient: the local Whittaker factor of the $N$-scaled binary lattice
at a split level prime, in closed form (Theorem~\ref{thm:closed}), with the corollary that the
empirical ``level support rule'' of the computational literature is exactly the vanishing of that
factor (Corollary~\ref{cor:support}). Section~\ref{sec:reenum} re-runs the case analysis with this
factor and finds no new case (Proposition~\ref{prop:reenum}). Section~\ref{sec:data} records the
correction that the CM values require, computed exactly, and Section~\ref{sec:elim} lists the
explanations we have eliminated. Section~\ref{sec:open} states what remains open.

All numerical statements below are exact rational computations, carried out in Magma; the
supporting scripts are described in Appendix~\ref{app:repro}.

\section{The chain, and the scope of its central step}\label{sec:chain}

Write $U=\lambda^\perp$ for the negative definite plane attached to a CM point of discriminant $d$,
$\Lp=L\cap\Q\lambda$ and $\Lm=L\cap U$. Following \cite[(11),(12)]{Yang}, one sets
\[
  \kappa_\eta(m)=\sum_{\mu\in L/(\Lp+\Lm)}\ \sum_{x\in \eta_++\mu_++\Lp}
     \kappa^-_{\eta_-+\mu_-}\!\bigl(m-\langle x,x\rangle/2\bigr),
\]
where $\kappa^-_\mu(m)$ is the $s$-linear coefficient of the $m$-th Fourier coefficient
$A_\mu(s,m,v)$ of an incoherent Eisenstein series of weight $1$ attached to $\Lm$, and
$\kappa^-_\mu(0)=0$ unless $\mu=0$.

For $m>0$ one has, with $S_{m,\mu}=\{p: p\mid \Delta \text{ or } v_p(m)>0\}$ and
$\Delta=\disc(\Lm)$,
\begin{equation}\label{eq:euler}
  A_\mu(s,m,v)=-\frac{2\pi}{L(s+1,\chi_d)}\prod_{p\in S_{m,\mu}}
    \frac{W_{m,p}(s,\varphi_{\mu,p})}{1-\chi_d(p)p^{-1-s}},
\end{equation}
a genuine Euler product because $\varphi_\mu=\mathrm{char}(\mu+\widehat{\Lm})$ is a pure tensor.
Since $A_\mu(s,m,v)$ vanishes at $s=0$, some factor must vanish there, and differentiating
\eqref{eq:euler} yields \cite[Lemma 16]{Yang}: the derivative is the single term at the vanishing
place, and vanishes outright if two or more factors vanish.

That step is correct as an identity. What it inherits from \cite{KRY} is the identification of the
vanishing of a local factor with a \emph{representability} obstruction, i.e.\ with membership of
$\Diff(m)$. Kudla--Rapoport--Yang state the scope of their computation plainly:

\begin{quote}
``in section 2, we have considered only those Eisenstein series built from sections which are
defined by the characteristic functions of the completions of $\mathcal{O}_k$. More general sections
could be considered at the cost of (i) more elaborate calculations of Whittaker functions and their
derivatives used in computing the Fourier expansion of the central derivative, and (ii)
incorporation of a level structure in the moduli problem.'' \cite[\S6]{KRY}
\end{quote}

\noindent
(They also fix a maximal order $\mathcal{O}_B$.) The lattice $\Lm$ arising here is not of that
shape. Write $\Delta=\disc(\Lm)$. For a discriminant $d$ with $N\nmid d$ --- the case in which the
level prime is \emph{unramified} in $k=\Q(\sqrt d)$, and hence, by Eichler's condition on optimal
embeddings, necessarily \emph{split} --- one has $\ord_N(\Delta)=2$: the binary lattice is
$N$-scaled at the level prime, not unimodular. This is exactly a ``more general section''.

\begin{remark}
The distinction is invisible in the maximal-order setting: there $\Lm$ is unimodular at every split
prime, no split prime can obstruct representability, and correspondingly the logarithms produced by
\cite[Theorem 4.1]{Schofer} run over ramified and inert places only. Every discriminant admitting a
CM point on $X_0^D(N)$ with $N\nmid d$ has $N$ split, so this configuration is precisely the one the
maximal-order theory never meets.
\end{remark}

\section{The local Whittaker factor at a split level prime}\label{sec:closed}

Part (i) of the programme quoted above, at the level prime, is the following computation. Let
$X=N^{-s}$.

\begin{theorem}\label{thm:closed}
Let $d<0$ be a discriminant with $N\nmid d$ and let $\Lm$ be the negative definite binary lattice
attached to a CM point of discriminant $d$, so $\ord_N(\disc \Lm)=2$. Then for every $m\ge 1$, with
$j=\ord_N(m)$,
\[
  W_{m,N}(X)\;=\;1+(N-1)\bigl(X+X^2+\cdots+X^{j}\bigr)-X^{\,j+1}.
\]
In particular
\[
  W_{m,N}(1)=(N-1)\,\ord_N(m),
  \qquad
  W_{m,N}'(1)=(j+1)\Bigl(\tfrac{(N-1)j}{2}-1\Bigr).
\]
\end{theorem}

The statement is a computation rather than a proof: it has been verified in exact arithmetic for
$N=2,3,5$ and all $1\le m\le 60$, on the lattices $\Lm$ of $X_0^{15}(2)$, $X_0^{6}(5)$ and
$X_0^{10}(3)$ at firing discriminants, giving $180$ independent checks. It is pinned as a regression
test in the accompanying code.

\begin{corollary}\label{cor:support}
$W_{m,N}(1)=0$ if and only if $N\nmid m$.
\end{corollary}

This identifies, and explains, an empirical rule that has been observed repeatedly in computations
with these curves: the correction discussed in Section~\ref{sec:data} is supported exactly on the
indices $m$ with $N\mid m$. That rule is the vanishing of the level-prime factor.

\begin{remark}
At $j=0$ the factor is $W_{m,N}(X)=1-X=1-N^{-s}$, the inverse local zeta factor. Its vanishing at
$s=0$ is therefore of a different nature from a representability obstruction: it holds identically
for all $m$ with $N\nmid m$, and at a split $N$ the local quadratic space represents every $m$.
\end{remark}

\begin{remark}
The value $W_{m,N}(1)=(N-1)\ord_N(m)$ coincides with the local factor Kudla--Yang obtain at $p\mid N$
for the \emph{quaternary} Eichler lattice \cite[(9.1)]{KY}. The agreement appears to be a genuine
coincidence of shape between the quaternary lattice and the binary $\Lm$; we do not have an
explanation for it.
\end{remark}

\section{Re-running the case analysis}\label{sec:reenum}

With Theorem~\ref{thm:closed} in hand one can redo the enumeration of \cite[Prop.~2.8]{KRY} for the
$N$-scaled section. Normalising as in \eqref{eq:euler}, put
$f_p(s)=W_{m,p}(s)/(1-\chi_d(p)p^{-1-s})$.

\begin{proposition}\label{prop:reenum}
At a firing discriminant, $f_N$ has a simple zero at $s=0$ exactly when $N\nmid m$, and this zero is
not a $\Diff$ zero. Consequently
\[
  \ord_{s=0}A_\mu(s,m,v)\;=\;\lvert\Diff(m)\rvert+[\,N\nmid m\,],
\]
and since incoherence forces $\lvert\Diff(m)\rvert\ge 1$:
\begin{enumerate}
\item if $N\mid m$, the order is $\lvert\Diff(m)\rvert$, and the derivative is nonzero precisely when
      $\lvert\Diff(m)\rvert=1$, contributing a logarithm at the ramified or inert place of $\Diff(m)$;
\item if $N\nmid m$, the order is at least $2$ and the derivative vanishes.
\end{enumerate}
In neither case does a $\log N$ occur.
\end{proposition}

So the level structure only ever raises the order of vanishing: it removes contributions, and
creates no new contributing configuration. The three cases of \cite[Prop.~2.8]{KRY} stand, with no
fourth. This is confirmed computationally: over roughly $190$ terms $(\mu,x)$ arising in the
evaluation of $\kappa_0(m)$ on three bases, the level prime is \emph{never} the unique vanishing
place, while it lies among two or more in $60$--$80\%$ of terms.

\begin{corollary}\label{cor:onlyplace}
Within this framework, a $\log N$ can enter a CM value only through the $m=0$ term
$c_0(0)\kappa_0(0)$, where $\kappa_0(0)$ carries the summand $\sum_{p\mid N/(N,d)}\log p$ of
\cite[Remark 15(2), Lemma 20]{Yang}.
\end{corollary}

\section{The correction the CM values require}\label{sec:data}

Against this stands a computation of the input form itself. Following \cite[eq.~(6)]{GY} one may
form
$F_f=\sum_{\gamma\in\Gamma_0(M)\backslash SL_2(\Z)}(f\vert_{1/2}\gamma)\,\rho(\gamma^{-1})e_0$
and evaluate it numerically; the metaplectic phase is handled by evaluating both the slash and
$\rho$ along one and the same $ST$-word, so that they share a single metaplectic lift. Every such
evaluation is gated on reproducing the principal part of $F_f$ predicted by \cite[Lemma 24]{GY} from
the two scalar $q$-expansions of $f$; on the runs reported here that gate is met to $10^{-103}$.

\begin{observation}\label{obs:c00}
For every Borcherds form produced by the pipeline on $X_0^{15}(2)$, $X_0^{6}(5)$, $X_0^{10}(3)$ and
$X_0^{21}(2)$ one finds $c_0(0)=0$ (to $10^{-103}$), so the hypothesis $c_0(0)=0$ of
\cite[Thm.~B]{GY} is satisfied and, by Corollary~\ref{cor:onlyplace}, no $\log N$ can arise.
\end{observation}

\begin{observation}\label{obs:iso}
The isotropic cosets of $\Ldual/L$ number exactly $2N-1$; the $D$-part of the discriminant form is
anisotropic, so all of them come from the level-$N$ hyperbolic plane. All $2N-2$ nonzero isotropic
cosets carry the same constant term $c_\eta(0)$.
\end{observation}

\begin{observation}\label{obs:mult}
Let $\mathrm{mult}(f)=\tfrac12 c_\eta(0)$ for any nonzero isotropic $\eta$. Then adding
$\mathrm{mult}(f)\cdot\sum_{p\mid N/(N,d_{\mathrm{fund}})}\log p$ to the computed CM value
reproduces, exactly, the corrections independently measured by kernel-consistency on
$X_0^{15}(2)$ ($9$ values), $X_0^{6}(5)$ ($5$) and $X_0^{10}(3)$ ($5$), and by the consistency
relation of Observation~\ref{obs:haupt} below on $X_0^{14}(5)$ ($1$, at eleven digits): $20$ of
$20$. A further base $X_0^{21}(2)$ supplies nine values with no prior ground truth against which
to check them.
\end{observation}

That the correction is genuinely required, and not an artefact of a normalisation, is visible on
$X_0^{6}(19)$: on a set of CM points including discriminants divisible by $N$, the uncorrected
values are mutually \emph{inconsistent} --- the associated system has no solution --- whereas with
the correction the published model is recovered.

\subsection*{Ground truth without an Eisenstein computation}
The pipeline imposes, at every rational CM point $d$, the consistency relation
\begin{equation}\label{eq:haupt}
  \varepsilon_1(d)\,\frac{s(d)}{s(d_A)}+\varepsilon_2(d)\,\frac{\tilde s(d)}{\tilde s(d_B)}=1,
  \qquad \varepsilon_i(d)\in\{\pm1\},
\end{equation}
between the two Hauptmodul rows, normalised at the discriminants $d_A$, $d_B$ where they vanish.
Restoring a correction $r_f\log N$ to $\log\lvert\psi_f\rvert$ multiplies the normalised value at
$d$ by $N^{\,r_f(k(d)-k(d_{\mathrm{ref}}))}$, where $k(d)\in\{0,1\}$ records whether the term
fires at $d$. Each rational CM point whose firing status differs from the reference's therefore
contributes one equation in $(N^{r_1},N^{r_2})$; two determine the pair and any further points are
consistency checks. This measures the multipliers of the two Hauptmodul forms on any base the
pipeline reaches, in seconds, with no Eisenstein series evaluated.

\begin{observation}\label{obs:haupt}
On nine bases the relation \eqref{eq:haupt} determines the pair $(r_{-1},r_{-2})$ uniquely:
\[
\begin{array}{llll}
X_0^{15}(2)\colon (4,2) & X_0^{14}(5)\colon (1,1) & X_0^{6}(7)\colon (1,2) & X_0^{6}(5)\colon (0,0)\\
X_0^{10}(3)\colon (0,0) & X_0^{10}(11)\colon (0,0) & X_0^{6}(19)\colon (0,0) & X_0^{22}(7)\colon (0,0)\\
X_0^{34}(5)\colon (0,0) & & &
\end{array}
\]
Every value with an independent determination is reproduced --- $X_0^{15}(2)$ against the
kernel-consistency measurement, $X_0^{6}(5)$ and $X_0^{10}(3)$ against the numerical evaluation of
$F_f$, and $X_0^{6}(7)$, the one \emph{nonzero} cross-check, against an earlier independent
measurement. Conversely, on $X_0^{14}(5)$ the pair $(1,1)$ was measured first by
\eqref{eq:haupt} and \emph{then} confirmed by the numerical evaluation of $F_f$:
$\tfrac12 c_\eta(0)=1.00000000000$ on the form with pole order $16$, with the principal-part gate
met to $2.3\cdot10^{-37}$ and all eight nonzero isotropic cosets agreeing. This is the first
verification of the rule of Observation~\ref{obs:mult} at $N>2$ against ground truth obtained by
an entirely independent method.
\end{observation}

Writing $\mathrm{mult}(f)=\sum_m c(-m)A_m$ and using Borcherds' obstruction to set
$A_m=-b(m)/4$ with $b$ the relevant weight $3/2$ Eisenstein coefficient, the linear functional is
determined per base by the measured multipliers. Without further hypotheses the system leaves four
free directions on each base; imposing the support rule of Corollary~\ref{cor:support} (the weight
vanishes unless $N\mid m$) removes them all, and the resulting weights are \emph{unique} with six
spare conditions satisfied on each of the three bases --- eighteen independent chances to falsify
the support rule, none taken. Table~\ref{tab:A} records the result.

\begin{table}[ht]
\centering
\begin{tabular}{llll}
\toprule
base & \multicolumn{3}{l}{$A_m$ (equivalently $b(m)=-4A_m$)}\\
\midrule
$X_0^{15}(2)$ & $A_2=-1$   & $A_{10}=1$   & $A_{30}=0$\\
$X_0^{6}(5)$  & $A_{10}=3/2$ & $A_{15}=1/2$ & $A_{30}=3/2$\\
$X_0^{10}(3)$ & $A_3=1/2$  & $A_{12}=1/2$ & $A_{30}=3/2$\\
$X_0^{21}(2)$ & $A_2=-2$, $A_6=-4$ & $A_{18}=2$, $A_{42}=-8$ & ($0$-block: $B_{21}=-5$)\\
\bottomrule
\end{tabular}
\medskip
\caption{The correction, per index, as a linear functional of the principal part. The
values are $d$-independent: at every firing discriminant the coefficient of $\log N$ in
$\kappa_0(m)$ must equal $A_m$, whereas the recipe of \S\ref{sec:chain} returns $0$.}
\label{tab:A}
\end{table}

\section{Explanations eliminated}\label{sec:elim}

Each of the following was tested against the data of Section~\ref{sec:data} and is inconsistent
with it. We record them because several are natural first guesses.

\begin{enumerate}
\item $\mathrm{mult}=c_0(0)$, or any linear combination of the two scalar constant terms
      $\mathrm{Coeff}(f_\infty,0)$ and $\mathrm{Coeff}(f_0,0)$: inconsistent on three bases.
\item Evaluating the local Whittaker at a nonzero isotropic coset $\eta^*$ rather than at $0$:
      destroys the pattern of vanishing indices on all three bases. The mechanism is now clear:
      the nonzero isotropic cosets live in the level-$N$ hyperbolic plane, so changing the coset
      changes \emph{only} the level-prime factor (measured: from $1-X^2/N$ to the constant $1$,
      with the factors at $p\mid D$ byte-identical) and can never repair an index whose defect is
      at a ramified prime.
\item Replacing the local factor at $p\mid D$ by the anisotropic branch of \cite[Prop.~5.3]{KY}:
      does not reproduce $b$, under either sign of $\kappa m$. (The implementation of
      \cite[Props.~5.3, 5.4]{KY} agrees with the pipeline's own factor at every good prime,
      $470$ checks, and at the level prime, $150$ checks.)
\item The incoherent support condition $\lvert\Diff(m)\rvert=1$: scores $8/9$ on the exact
      coefficients and is the only rule that predicts $b(30)=0$ on $X_0^{15}(2)$, but is far too
      permissive --- for $r\le 60$ it holds at $15$, $19$ and $20$ indices with $N\nmid r$, where
      the coefficient must vanish.
\item Cancelling the zeta factor $1-X$ at the level prime before counting vanishing places: the
      recovered terms differentiate at ramified and inert primes, never at $N$.
\item The inner $m=0$ of $\kappa_\eta$, i.e.\ terms with $\langle x,x\rangle/2=m$: never occurs, in
      all $60$ (discriminant, index) pairs over every firing discriminant of the three bases.
\item Schofer's outer $m=0$ term: vanishes by Observation~\ref{obs:c00}.
\item The re-enumeration of \S\ref{sec:reenum}: adds no case.
\item Any re-weighting of the existing local data: every combination of a subset of the
      prefactors $\{\sqrt{m}/\sqrt{\lvert d\rvert},\,h/w,\,\prod(1-\chi(p)/p),\,
      \prod(1-1/p^2)^{-1}\}$ with a local rule drawn from \{the product of the $W_p(1)$; the
      derivative of that product; the derivative taken at $N$, at the primes $p\mid D$, at each
      single prime, or at the vanishing places\}, times an arbitrary constant per base, fails to
      reproduce the nine measured multipliers of $X_0^{15}(2)$.
\item More strongly, \emph{no multiplicative correction of any kind} can succeed: at two of the
      nine determined indices (both with $m=10$) the product of the code's local factors vanishes
      while the true weight does not, so the discrepancy is infinite there. The missing
      contribution is a \emph{term}, not a factor.
\item Differentiating at the place where the ternary \emph{space} fails to represent $m$
      (Kudla's $\Diff$, computed at the level of quadratic spaces rather than lattice densities):
      reproduces $A_2=-1$ and $A_{30}=0$ on $X_0^{15}(2)$ exactly --- the only rule found that
      \emph{explains} $A_{30}=0$ --- but fails on $7$ of $9$ indices, and is refuted as a support
      rule by $X_0^{6}(5)$ at $m=15$, where $\Diff$ is empty but $A_{15}=1/2\ne0$. The support is
      $N\mid m$; whether $\Diff$ supplies the correct place to differentiate is a separate
      question.
\end{enumerate}

\section{The correction identified}\label{sec:ident}

The rule of Observation~\ref{obs:mult} reads off the constant term $c_{\eta^*}(0)$ of the input
form at a nonzero isotropic coset $\eta^*$. By the constant-term pairing (pair $F$ against a
holomorphic Eisenstein series of the complementary weight $3/2$ and take the constant term of the
resulting weight-$2$ invariant), that quantity is determined by the principal part of $F$ paired
against the coefficients of $E_{A,\eta^*}$, the Eisenstein series attached to the coset $\eta^*$
of the discriminant form $A=\Ldual/L$. This is precisely the object the maximal-order literature
does not compute: \cite{BK}, \cite[\S8--9]{KY} and the recipe of \S\ref{sec:chain} all work with
the series attached to the \emph{zero} coset (equivalently, with sections
$\mathrm{char}(\mu+\widehat L)$ seeded at $0$), and Kudla--Yang explicitly leave the general
Eichler-order case open at the end of their \S9.

The coefficients of $E_{A,\beta}$ for isotropic $\beta\ne0$ are computed by Schwagenscheidt
\cite{Schwag}. For our purposes the essential ingredient is not the main formula but the oldform
relation \cite[Prop.~1.4]{Schwag} applied to the trivial character: for $\beta=\eta^*$ of prime
order $N_\beta=N$,
\begin{equation}\label{eq:oldform}
  E_{A,\eta^*,\chi_0}\;=\;E_{B,0}\!\uparrow_B^A\;-\;E_{A,0},
  \qquad
  B=\langle\eta^*\rangle^{\perp}/\langle\eta^*\rangle\;\cong\;L_1^{\vee}/L_1,
\end{equation}
where $L_1$ is the even lattice generated by $L$ and $\eta^*$. Pairing $F$ against
\eqref{eq:oldform} and using Observations~\ref{obs:c00} and~\ref{obs:iso} --- $c_0(0)=0$, the
constant $c_\eta(0)$ shared by all $2N-2$ nonzero isotropic cosets, and the vanishing of the
constant-term pairing for every nontrivial character twist --- gives
\begin{equation}\label{eq:multformula}
  \mathrm{mult}(f)\;=\;-\frac{1}{4(N-1)}\ \sum_{n>0}\ \sum_{\gamma\in(\eta^*)^{\perp}}
     c_\gamma(-n)\,b^{B,0}_{\gamma+H}(n),
  \qquad H=\langle\eta^*\rangle,
\end{equation}
in which every object on the right is a \emph{seed-$0$} Eisenstein coefficient of the smaller
module $B$ --- the territory of \cite{BK}, computable by the same machinery as the existing
recipe, applied to $L_1$ in place of $L$.

Three structural facts corroborate the identification.
\begin{enumerate}
\item The pairing constant in \eqref{eq:multformula} is $4(N-1)$ --- exactly the normalisation
      measured in Observation~\ref{obs:mult} (the sum of $c_\eta(0)$ over the $2N-2$ nonzero
      isotropic cosets equals $4(N-1)\cdot\mathrm{mult}$).
\item $\det L_1=\det L/N^2$: the level-prime part of $L_1$ is \emph{unimodular}, so the
      coefficients $b^{B,0}$ have no vanishing level-prime factor. The zero
      $W_{m,N}(1)=0$ at $N\nmid m$ (Corollary~\ref{cor:support}), which removes the majority of
      terms from the seed-$0$ pairing on $A$, never occurs on $B$ --- the structural mechanism by
      which \eqref{eq:multformula} can produce the $\log N$ terms that \S\ref{sec:reenum} shows
      the seed-$0$ recipe cannot.
\item $2\lvert\det L_1\rvert$ remains a perfect square (e.g.\ $2\cdot450=30^2$ on
      $X_0^{15}(2)$), so the discriminant and character bookkeeping of the weight-$3/2$
      coefficient formula carries over unchanged in shape.
\end{enumerate}

Two cautions. The identity \eqref{eq:oldform} is proven in \cite{Schwag} for weight
$k\ge5/2$; at $k=3/2$ both sides require analytic continuation, and the pairing argument assumes
the continued series remains holomorphic at the special point, where Zagier-type non-holomorphic
contributions are possible in principle. And at rank $3$ and weight $3/2$ the local factors of
$b^{B,0}$ collapse to plain representation densities of $L_1$ (the Bruinier--Kuss local
polynomial is evaluated where its factor $1-p^{m-1}X$ vanishes, leaving only the stabilised
leading term), so the formula is elementary to evaluate; its direct numerical confrontation with
the nine measured multipliers of $X_0^{15}(2)$ is the immediate next step, in progress at the
time of writing.

\section{The open question}\label{sec:open}

Corollary~\ref{cor:onlyplace} and Observation~\ref{obs:c00} are in direct tension with
Observation~\ref{obs:mult} and with the failure of the uncorrected computation on $X_0^{6}(19)$. We
see two possibilities and cannot presently distinguish them.

\begin{enumerate}
\item The constant term $c_0(0)$ of the form $F$ actually attached to $\psi_F$ differs from that of
      the coset sum $F_f$. The principal parts agree to $10^{-103}$, but the principal part alone
      does not determine $c_0(0)$.
\item The correction compensates an error elsewhere in the computation of the CM values, outside
      Schofer's formula. It is a $d$-independent linear functional of the principal part, which is
      the shape of a per-form normalisation.
\end{enumerate}

The identification of \S\ref{sec:ident} sharpens the second reading. The $m=0$ term of the
formula, as transmitted through \cite{Schofer,Yang}, collapses the sum $\sum_\eta c_\eta(0)
\kappa_\eta(0)$ to the single term $c_0(0)\kappa_0(0)$ via a condition of the form
$\delta_{0,\mu}$. But the collapse condition constrains only the \emph{negative-definite} part of
the coset, so a nonzero isotropic $\eta$ with $\eta_-=0$ is not excluded by it --- and by
Observations~\ref{obs:iso} and~\ref{obs:mult} it is exactly those cosets, with their shared
nonzero constant term, that carry the correction. On this reading nothing outside Schofer's
formula is being compensated; rather, terms inside it were dropped.

In either case Theorem~\ref{thm:closed} and Proposition~\ref{prop:reenum} stand, and part (ii) of
the Kudla--Rapoport--Yang programme --- the incorporation of a level structure --- remains to be
carried out. Table~\ref{tab:A} is what any corrected formula must reproduce, index by index.

A third reading should be recorded, because it is not excluded by anything above. The evaluation of
$\kappa^-_\mu(m)$ for $m>0$ has never been checked against an independent computation at $N>1$ ---
neither here nor, as far as we can determine, in the literature. Yang's worked examples
\cite[Ex.~21--24]{Yang} take $B=\bigl(\tfrac{-1,3}{\Q}\bigr)$ with a \emph{maximal} order, and
indeed $L=\Lp+\Lm$ there; the corresponding regression tests in the accompanying code likewise cover
only $N=1$. What Observation~\ref{obs:haupt} now supplies is a verification at $N>1$ of the
\emph{correction rule} --- the multiplier $\tfrac12 c_\eta(0)$, confirmed at $N=5$ against ground
truth measured by an independent method --- but the evaluation of $\kappa^-_\mu(m)$ itself, which
sits upstream in the $m>0$ terms, remains unverified at $N>1$.

\begin{remark}[Published CM values need not discriminate]\label{rem:gy36}
A natural-seeming check is to compare against a published CM value on a base with $N>1$ and a
nonvanishing correction; $X_0^{14}(5)$ with $s(\tau_{-280})=5/16$ \cite[Ex.~36]{GY} appears to be
the only available instance. It cannot work. The Hauptmodul there is normalised by
$s(\tau_{-4})=\infty$, $s(\tau_{-11})=1$, $s(\tau_{-35})=0$, so the published value depends only on
the ratio $s(-280)/s(-35)$ --- and $-35$ and $-280$ are both non-firing ($5\mid d$), so any
correction multiplies both by the same power of $N$ and cancels; the remaining anchors sit at $0$
and $\infty$, where a multiplicative factor is invisible. Indeed the pipeline reproduces $5/16$
exactly both with and without the correction. A published value can only test the correction if
its normalising anchors and the tested point lie on opposite sides of the firing rule; this should
be checked, at a cost of nothing, before any such comparison is attempted. The consistency
relation \eqref{eq:haupt}, by contrast, is sensitive exactly where \cite[Ex.~36]{GY} is blind, and
is what supplied the $X_0^{14}(5)$ ground truth here.
\end{remark}

\appendix
\section{Reproducibility}\label{app:repro}

The computations were carried out in Magma against the codebase accompanying this note.

\begin{itemize}
\item \texttt{VectorValuedForm.m} constructs $F_f$: coset representatives for
      $\Gamma_0(M)\backslash SL_2(\Z)$, $ST$-words, the weight $1/2$ metaplectic slash, the Weil
      representation over $\C$, and the Fourier extraction. The action of $\rho(S)$ is implemented
      as a finite Fourier transform on $\Ldual/L$, which reduces the cost from $\lvert G\rvert^2$ to
      $\lvert G\rvert\sum_r d_r$ and the memory from $\lvert G\rvert^2$ to $\lvert G\rvert$; without
      this the case $\lvert G\rvert=24200$ is out of reach.
\item \texttt{EisensteinLocalFactors.m} implements \cite[(2.16),(2.17), Props.~5.3, 5.4]{KY} and
      exposes \texttt{LocalWhittakerPolynomial}, the Whittaker polynomial whose $s$-derivative
      produces the logarithms.
\item Numerical parameters matter. The horocycle height must satisfy $\Im\tau\ge1$, or some coset
      is carried into the pole at $\infty$ and the terms overwhelm their sum. The number of sample
      points must beat the positive Fourier coefficients, which grow like $\exp(4\pi\sqrt{pm})$ for
      pole order $p$: on the pole-order-$30$ forms of $X_0^{15}(2)$ the principal-part gate reads
      $5.8\cdot10^{73}$, $1.2\cdot10^{62}$, $2.7\cdot10^{28}$ and $8.7\cdot10^{-103}$ at $K=48$,
      $64$, $96$ and $192$ respectively. Raising the working precision without raising $K$ has no
      effect.
\item Every reported value is gated on the principal-part check; a run failing that gate is
      discarded.
\end{itemize}

\begin{thebibliography}{99}

\bibitem{BK} J.~H.~Bruinier and M.~Kuss, \emph{Eisenstein series attached to lattices and modular
forms on orthogonal groups}, Manuscripta Math.\ \textbf{106} (2001), 443--459.

\bibitem{Errthum} E.~Errthum, \emph{Singular moduli of Shimura curves},
Canad.\ J.\ Math.\ \textbf{63} (2011), 826--861.

\bibitem{GY} J.-W.~Guo and Y.~Yang, \emph{Equations of hyperelliptic Shimura curves},
Proc.\ Lond.\ Math.\ Soc.\ (2017), doi:10.1112/S0010437X16007739.

\bibitem{KRY} S.~Kudla, M.~Rapoport and T.~Yang, \emph{On the derivative of an Eisenstein series of
weight one}, Internat.\ Math.\ Res.\ Notices (1999), no.~7, 347--385.

\bibitem{KY} S.~Kudla and T.~Yang, \emph{Eisenstein series for $SL(2)$},
Sci.\ China Math.\ \textbf{53} (2010), 2275--2316.

\bibitem{Schofer} J.~Schofer, \emph{Borcherds forms and generalizations of singular moduli},
J.\ reine angew.\ Math.\ \textbf{629} (2009), 1--36.

\bibitem{Schwag} M.~Schwagenscheidt, \emph{Eisenstein series for the Weil representation},
J.\ Number Theory \textbf{193} (2018), 74--101.

\bibitem{Yang} Y.~Yang, \emph{Special values of hypergeometric functions and periods of CM elliptic
curves}, arXiv:1503.07971.

\end{thebibliography}

\end{document}
